Because \(k_b=M+1\), the convention in \texttt{def:rank-inversion} gives \(Q_{k_b}=+\infty\). Hence, for every realization and every boundary-arm score map,
\[
s_b(O_{T,b})\le Q_{k_b}(s_b(O_{1,b}),\ldots,s_b(O_{M,b}))=+\infty.
\tag{1}
\]
Thus the boundary-arm event is the entire sample space. Removing it from the conjunction in \texttt{def:coverage-frontier} proves that the boundary-arm score map, preorder, target lift, and calibration-orbit coordinate are irrelevant.

Fix \(s_c\) and condition on \(O_{T,c}=O_t\). The elementary order-statistic identity gives
\[
s_c(O_t)\le Q_{k_c}(s_c(O_{1,c}),\ldots,s_c(O_{M,c}))
\quad\Longleftrightarrow\quad
N_t:=\sum_{i=1}^M\mathbf 1\{s_c(O_{i,c})<s_c(O_t)\}<k_c.
\tag{2}
\]
Indeed, the \(k_c\)-th order statistic is strictly below the target score exactly when at least \(k_c\) calibration scores are strictly below it. By \texttt{ass:known-matched-randomization}, \texttt{ass:selector-auxiliary-randomization}, and \texttt{def:observable-selector-mixture}, the nonboundary calibration orbits are iid with marginal mass function \(q_c\). They are independent of the target by \texttt{ass:prespecified-target}. Therefore, conditional on \(O_{T,c}=O_t\), the summands in \(N_t\) are iid Bernoulli variables with success probability
\[
\theta(O_t)=\sum_{O\in\mathcal O_c:s_c(O)<s_c(O_t)}q_{c,O}.
\tag{3}
\]
Consequently, \(N_t\sim\operatorname{Bin}(M,\theta(O_t))\). Averaging (2) over the target-orbit law proves the exact formula
\[
\sum_{O_t\in\mathcal O_c}p_{c,O_t}
\Pr\!\left\{\operatorname{Bin}\bigl(M,\theta(O_t)\bigr)<k_c\right\}.
\tag{4}
\]
The strict inequality in (3) shows explicitly that (4) handles ties.

It remains to sharpen the evaluator. Fix one total preorder and write its ordered nonempty tie classes as
\[
C_1\prec C_2\prec\cdots\prec C_d,\qquad d\le m_c.
\]
Set
\[
q_j=\sum_{O\in C_j}q_{c,O},\qquad p_j=\sum_{O\in C_j}p_{c,O},\qquad \theta_j=\sum_{\ell<j}q_\ell.
\tag{5}
\]
Every target orbit in \(C_j\) has the same value \(\theta_j\), so (4) becomes
\[
\sum_{j=1}^d p_j\Pr\!\left\{\operatorname{Bin}(M,\theta_j)<k_c\right\}.
\tag{6}
\]
The quantities in (5) and the sum in (6) are computed in one ordered pass: initialize \(\theta_1=0\), scan the members of \(C_j\) once to form \(q_j\) and \(p_j\), evaluate the current binomial probability, and update \(\theta_{j+1}=\theta_j+q_j\). Each of the \(m_c\) orbits is visited exactly once.

By \texttt{lem:exact-binomial-shorter-tail-evaluator}, for \(0<\theta_j<1\) the binomial probability in (6) is evaluated exactly in
\[
O\!\left(\log M+\min\{k_c,M-k_c+1\}\right)
\tag{7}
\]
arithmetic operations and constant auxiliary memory. When \(\theta_j=0\), the probability equals one because \(k_c\ge1\); when \(\theta_j=1\), it equals zero because \(k_c\le M\). These endpoints require neither division nor a recurrence. Thus one preorder costs
\[
O\!\left(m_c\left[1+\log M+\min\{k_c,M-k_c+1\}\right]\right)
\tag{8}
\]
arithmetic operations and constant auxiliary CDF working memory after the tabulated masses and the current preorder representation are available. The rank event depends on \(s_c\) only through its total preorder, and every total preorder has a representative in \([0,K]\). Enumerating all \(F_{m_c}\) preorders therefore gives the sharper total bound
\[
O\!\left(F_{m_c}m_c\left[1+\log M+\min\{k_c,M-k_c+1\}\right]\right).
\tag{9}
\]
Since \(m_c\ge1\), \(k_c\ge1\), \(\log M=O(M)\), and \(\min\{k_c,M-k_c+1\}\le k_c\), (9) also implies the looser displayed bound in the frozen statement, while constant auxiliary memory implies its \(O(k_c)\) bound. Formula (6) uses only the nonboundary marginals, so no factor \(F_{m_{1-c}}\), member-level sum \(|\mathcal V|\), or per-preorder \(m_0m_1\) scan occurs.